\uuid{iYFi}
\exo7id{3064}
\titre{exo7 3064}
\auteur{quercia}
\organisation{exo7}
\datecreate{2010-03-08}
\isIndication{false}
\isCorrection{false}
\chapitre{Propriétés de R}
\sousChapitre{Autre}
\cpge{mpsi-mp2i.nombres-reels-suites.a, pcsi.raisonnement-vocabulaire-ensembliste.c, ptsi.raisonnement-vocabulaire-ensembliste.c, mpsi-mp2i.nombres-reels-suites.g, mpsi-mp2i.structures-algebriques.b, bcpst1.nombres-reels, mp-mpi.structures-algebriques.a}
\module{Analyse}
\niveau{L1}
\difficulte{}

\contenu{
\texte{
Soit $H$ un sous-groupe additif de $\R$, $H \ne \{0\}$.
On pose $H^{+*} = H\cap\R^{+*}$, et $\alpha = \inf(H^{+*})$.
}
\begin{enumerate}
    \item \question{Si $\alpha \in H^{+*}$, montrer que $H = \alpha \Z$.}
    \item \question{Si $\alpha \notin H^{+*}$, montrer que $\alpha = 0$ et en d{\'e}duire que $H$
    est dense dans $\R$.}
\end{enumerate}
}
